

\chapter{Algebraic Number Theory Preliminaries}
The purpose of this appendix is to provide the algebraic number theory tools needed for the proof of Brauer's characterization of modular characters.

\section{Algebraic Integers}

\begin{defn}
A complex number \(\alpha \in \mathbb{C}\) is called an \textbf{algebraic integer} if there exist integers \(a_0, \dots, a_{n-1} \in \mathbb{Z}\) such that
\[
\alpha^n + a_{n-1}\alpha^{n-1} + \cdots + a_0 = 0.
\]
The set of all algebraic integers is denoted by \(\mathcal{R}\).
\end{defn}

\begin{prop}
\(\mathcal{R}\) is a ring. Moreover, if \(\alpha \in \mathcal{R}\), then \(\mathbb{Z}[\alpha]\) is a finitely generated \(\mathbb{Z}\)-module.
\end{prop}

\begin{proof}
If \(\alpha\) satisfies \(\alpha^n + a_{n-1}\alpha^{n-1} + \cdots + a_0 = 0\), then \(\alpha^n = -a_{n-1}\alpha^{n-1} - \cdots - a_0\). By induction, every power \(\alpha^k\) for \(k \ge n\) is a \(\mathbb{Z}\)-linear combination of \(1, \alpha, \dots, \alpha^{n-1}\). Thus \(\mathbb{Z}[\alpha]\) is generated by these \(n\) elements.

If \(\alpha, \beta \in \mathcal{R}\), then \(\mathbb{Z}[\alpha]\) and \(\mathbb{Z}[\beta]\) are finitely generated \(\mathbb{Z}\)-modules. The ring \(\mathbb{Z}[\alpha, \beta]\) is the image of the tensor product \(\mathbb{Z}[\alpha] \otimes_{\mathbb{Z}} \mathbb{Z}[\beta]\) under the multiplication map, hence is finitely generated over \(\mathbb{Z}\). Therefore every element of \(\mathbb{Z}[\alpha, \beta]\), including \(\alpha+\beta\) and \(\alpha\beta\), is integral over \(\mathbb{Z}\), hence lies in \(\mathcal{R}\).
\end{proof}

\begin{prop}\label{prop:integer}
If \(q \in \mathbb{Q} \cap \mathcal{R}\), then \(q \in \mathbb{Z}\).
\end{prop}

\begin{proof}
Write \(q = a/b\) in lowest terms with \(b > 0\). If \(q\) satisfies a monic polynomial over \(\mathbb{Z}\), then substituting \(q = a/b\) and multiplying by \(b^n\) gives
\[
a^n + b \cdot (\text{some integer}) = 0.
\]
Thus \(b \mid a^n\), so \(b = 1\). Hence \(q \in \mathbb{Z}\).
\end{proof}



\section{Trace and Discriminant}

\begin{defn}\label{def:ring of integers}
A \textbf{number field} is a finite extension \(K\) of \(\mathbb{Q}\). Its \textbf{ring of integers} is
\[
\mathcal{O}_K := K \cap \mathcal{R}.
\]
\end{defn}

For \(x \in K\), let \(\operatorname{Tr}_{K/\mathbb{Q}}(x)\) denote the trace of the \(\mathbb{Q}\)-linear map \(y \mapsto xy\).

\begin{lemma}
\label{lem:tr_int}
If \(x \in \mathcal{O}_K\), then \(\operatorname{Tr}_{K/\mathbb{Q}}(x) \in \mathbb{Z}\).
\end{lemma}

\begin{proof}
The trace is the sum of the Galois conjugates of \(x\). Each conjugate is an algebraic integer, so their sum is an algebraic integer. Since it is rational, it is an integer by Proposition\ref{prop:integer}.
\end{proof}


\begin{lemma}\label{lem:frac}
Every element of \(K\) can be written as \(a/b\) with \(a, b \in \mathcal{O}_K\), \(b \neq 0\).
\end{lemma}

\begin{proof}
Let \(x \in K\), and let its minimal polynomial over \(\mathbb{Q}\) be
\[
f(T) = T^n + c_{n-1}T^{n-1} + \cdots + c_0, \quad c_i \in \mathbb{Q}.
\]
Choose \(m \in \mathbb{Z}\) such that \(m c_i \in \mathbb{Z}\) for all \(i\). Then \(m x\) is an algebraic integer, so \(m x \in \mathcal{O}_K\). Hence \(x = (m x)/m\).
\end{proof}



\begin{lemma}
There exist \(\alpha_1, \dots, \alpha_n \in \mathcal{O}_K\) that form a \(\mathbb{Q}\)-basis of \(K\).
\end{lemma}

\begin{proof}
Take any \(\mathbb{Q}\)-basis \(\beta_1, \dots, \beta_n\) of \(K\). By Lemma \ref{lem:frac}, each \(\beta_i\) can be written as \(a_i/b_i\) with \(a_i, b_i \in \mathcal{O}_K\). Then \(\alpha_i = b_1 \cdots b_n \beta_i\) is in \(\mathcal{O}_K\), and the \(\alpha_i\)'s still form a \(\mathbb{Q}\)-basis.
\end{proof}

\begin{lemma}\label{lem:det}
Let \(\alpha_1, \dots, \alpha_n \in \mathcal{O}_K\) be a \(\mathbb{Q}\)-basis of \(K\). The discriminant
\[
\Delta := \det(\operatorname{Tr}(\alpha_i \alpha_j))_{i,j=1}^n
\]
is a non-zero integer.
\end{lemma}

\begin{proof}
The trace pairing \((x, y) \mapsto \operatorname{Tr}(xy)\) is non-degenerate because \(K/\mathbb{Q}\) is separable (characteristic 0). Hence the matrix is invertible over \(\mathbb{Q}\), so \(\Delta \neq 0\). By Lemma \ref{lem:tr_int}, all entries are integers, so \(\Delta \in \mathbb{Z}\).
\end{proof}

\begin{prop}\label{prop:finitely gen mod}
\(\mathcal{O}_K\) is a finitely generated \(\mathbb{Z}\)-module.
\end{prop}

\begin{proof}
Let \(\alpha_1, \dots, \alpha_n \in \mathcal{O}_K\) be a \(\mathbb{Q}\)-basis, and let \(\Delta\) be its discriminant. For any \(x \in \mathcal{O}_K\), write \(x = \sum_{i=1}^n q_i \alpha_i\) with \(q_i \in \mathbb{Q}\). Then
\[
\operatorname{Tr}(x \alpha_j) = \sum_{i=1}^n q_i \operatorname{Tr}(\alpha_i \alpha_j).
\]
The left-hand side is in \(\mathbb{Z}\) by Lemma \ref{lem:det}. Let \(M = (\operatorname{Tr}(\alpha_i \alpha_j))\). Then \(M \mathbf{q} = \mathbf{t}\) with \(\mathbf{t} \in \mathbb{Z}^n\). Since \(\Delta = \det M \neq 0\), we have
\[
\mathbf{q} = M^{-1} \mathbf{t} = \frac{1}{\Delta} \operatorname{adj}(M) \mathbf{t},
\]
so \(\Delta q_i \in \mathbb{Z}\) for all \(i\). Hence
\[
\mathcal{O}_K \subseteq \frac{1}{\Delta} \sum_{i=1}^n \mathbb{Z} \alpha_i,
\]
a finitely generated \(\mathbb{Z}\)-module. Therefore \(\mathcal{O}_K\) itself is finitely generated.
\end{proof}


\begin{lemma}[Determinant Lemma]
Let \(R\) be a commutative ring, and let \(M\) be an \(R\)-module generated by \(m_1, \dots, m_n\). Let \(\varphi: M \to M\) be an \(R\)-module homomorphism. Suppose there is an ideal \(I \subseteq R\) such that for each \(i\),
\[
\varphi(m_i) = \sum_{j=1}^n a_{ij} m_j, \quad a_{ij} \in I.
\]
Then \(\varphi\) satisfies a monic polynomial
\[
\varphi^n + b_{n-1}\varphi^{n-1} + \cdots + b_0 = 0,
\]
with \(b_k \in I\).
\end{lemma}

\begin{proof}
Let \(A = (a_{ij})\). Then for the column vector \(\mathbf{m} = (m_1, \dots, m_n)^T\),
\[
(\varphi I - A)\mathbf{m} = 0.
\]
Let \(B = \operatorname{adj}(\varphi I - A)\). Then
\[
\det(\varphi I - A) \mathbf{m} = B(\varphi I - A)\mathbf{m} = 0.
\]
Thus \(\det(\varphi I - A)\) annihilates each \(m_i\), hence annihilates all of \(M\). Therefore \(\det(\varphi I - A) = 0\) as an endomorphism of \(M\). Expanding \(\det(xI - A)\) gives a monic polynomial
\[
x^n + b_{n-1}x^{n-1} + \cdots + b_0,
\]
with each \(b_k\) a polynomial in the \(a_{ij}\)'s, hence \(b_k \in I\). Thus \(\varphi\) satisfies this polynomial.
\end{proof}

\begin{cor}
If \(x \in K\) satisfies \(x \mathcal{O}_K \subseteq \mathcal{O}_K\), then \(x \in \mathcal{O}_K\).
\end{cor}

\begin{proof}
Apply the determinant lemma with \(R = \mathbb{Z}\), \(M = \mathcal{O}_K\), \(\varphi\) = multiplication by \(x\), and \(I = \mathbb{Z}\). Since \(\varphi(\mathcal{O}_K) \subseteq \mathcal{O}_K\), the lemma gives a monic polynomial over \(\mathbb{Z}\) satisfied by \(x\). Hence \(x \in \mathcal{R} \cap K = \mathcal{O}_K\).
\end{proof}


\section{Prime Ideals and Maximal Ideals}

\begin{defn}
Let \(R\) be a commutative ring.
\begin{enumerate}
    \item An \textbf{ideal} is a subgroup \(I \subseteq R\) such that \(R I \subseteq I\).
    \item A \textbf{prime ideal} is a proper ideal \(\mathfrak{p}\) such that \(ab \in \mathfrak{p} \implies a \in \mathfrak{p}\) or \(b \in \mathfrak{p}\).
    \item A \textbf{maximal ideal} is a proper ideal \(\mathfrak{m}\) such that the only ideals containing \(\mathfrak{m}\) are \(\mathfrak{m}\) and \(R\).
\end{enumerate}
\end{defn}

\begin{prop}
In \(\mathcal{O}_K\), every non-zero prime ideal is maximal.
\end{prop}

\begin{proof}
Let \(\mathfrak{p} \neq 0\) be prime. Since \(\mathcal{O}_K\) is a finitely generated \(\mathbb{Z}\)-module (Proposition \ref{prop:finitely gen mod}), the quotient \(\mathcal{O}_K/\mathfrak{p}\) is finite. It is also an integral domain. A finite integral domain is a field. Hence \(\mathfrak{p}\) is maximal.
\end{proof}



\begin{defn}
For a non-zero ideal \(\mathfrak{p} \subseteq \mathcal{O}_K\), define
\[
\mathfrak{p}^{-1} := \{ x \in K \mid x \mathfrak{p} \subseteq \mathcal{O}_K \}.
\]
\end{defn}

\begin{lemma}\label{lem:inv}
For every non-zero prime ideal \(\mathfrak{p} \subseteq \mathcal{O}_K\), we have \(\mathfrak{p} \mathfrak{p}^{-1} = \mathcal{O}_K\).
\end{lemma}

\begin{proof}
Clearly \(\mathcal{O}_K \subseteq \mathfrak{p}^{-1}\), so \(\mathfrak{p} \subseteq \mathfrak{p} \mathfrak{p}^{-1} \subseteq \mathcal{O}_K\). Since \(\mathfrak{p}\) is maximal, either \(\mathfrak{p} \mathfrak{p}^{-1} = \mathfrak{p}\) or \(\mathfrak{p} \mathfrak{p}^{-1} = \mathcal{O}_K\).

We prove \(\mathfrak{p} \mathfrak{p}^{-1} \neq \mathfrak{p}\). Suppose otherwise. Then \(\mathfrak{p}^{-1} \subseteq \mathcal{O}_K\).

Choose \(0 \neq a \in \mathfrak{p}\). Let \(R = \mathcal{O}_K/(a)\), a finite ring. Let \(\bar{\mathfrak{p}}\) be the image of \(\mathfrak{p}\) in \(R\). Since \(R\) is finite, the descending chain
\[
\bar{\mathfrak{p}} \supseteq \bar{\mathfrak{p}}^2 \supseteq \bar{\mathfrak{p}}^3 \supseteq \cdots
\]
stabilizes: there exists \(n \ge 1\) such that \(\bar{\mathfrak{p}}^n = \bar{\mathfrak{p}}^{n+1}\).

We prove \(\bar{\mathfrak{p}}^n = 0\). Let \(M = \bar{\mathfrak{p}}^n\). Since \(R\) is finite, \(M\) is finite, hence finitely generated as an \(R\)-module. Let \(m_1, \dots, m_s\) be a generating set of \(M\) (for example, take all elements of \(M\)). Since \(M = \bar{\mathfrak{p}} M\), for each \(i\) there exist \(a_{ij} \in \bar{\mathfrak{p}}\) such that
\[
m_i = \sum_{j=1}^s a_{ij} m_j.
\]
Let \(A = (a_{ij})\). Then \((I - A)\mathbf{m} = 0\). Multiplying by \(\operatorname{adj}(I - A)\) gives
\[
\det(I - A) m_i = 0 \quad \text{for all } i.
\]
Since \(a_{ij} \in \bar{\mathfrak{p}}\), we have \(\det(I - A) \equiv 1 \pmod{\bar{\mathfrak{p}}}\). In the finite ring \(R\), this implies \(\det(I - A)\) is a unit. Hence \(m_i = 0\) for all \(i\), so \(M = 0\). Thus \(\bar{\mathfrak{p}}^n = 0\).

Therefore \(\mathfrak{p}^n \subseteq (a)\). Choose \(n\) minimal. Then there exists \(b \in \mathfrak{p}^{n-1} \setminus (a)\). Now
\[
\frac{b}{a} \cdot \mathfrak{p} \subseteq \frac{1}{a} \mathfrak{p}^n \subseteq \mathcal{O}_K,
\]
so \(b/a \in \mathfrak{p}^{-1}\). By our assumption \(\mathfrak{p}^{-1} \subseteq \mathcal{O}_K\), hence \(b/a \in \mathcal{O}_K\), so \(b \in (a)\), a contradiction. Thus \(\mathfrak{p} \mathfrak{p}^{-1} \neq \mathfrak{p}\), so \(\mathfrak{p} \mathfrak{p}^{-1} = \mathcal{O}_K\).
\end{proof}


\section{Unique Factorization of Ideals}

\begin{lemma}\label{product of non-zero prime ideals}
Every non-zero ideal \(I \subseteq \mathcal{O}_K\) contains a product of non-zero prime ideals.
\end{lemma}

\begin{proof}
Suppose not. Let \(S\) be the set of ideals that do NOT contain a product of prime ideals. Since \(\mathcal{O}_K\) is Noetherian (it is a finitely generated \(\mathbb{Z}\)-module), choose a maximal element \(I \in S\). Then \(I\) is not prime (otherwise it would contain itself, a product of one prime). So there exist \(a, b \notin I\) with \(ab \in I\). Let \(I_1 = I + (a)\) and \(I_2 = I + (b)\). Both strictly contain \(I\), so by maximality, neither is in \(S\). Hence each contains a product of prime ideals. Their product is contained in \(I_1 I_2 \subseteq I\). Thus \(I\) contains a product of prime ideals, contradiction. Therefore \(S = \emptyset\).
\end{proof}

\begin{thm}\label{thm:to prime ideals}
Every non-zero ideal \(I \subseteq \mathcal{O}_K\) can be written uniquely as
\[
I = \mathfrak{p}_1^{e_1} \cdots \mathfrak{p}_r^{e_r},
\]
where the \(\mathfrak{p}_i\)'s are distinct prime ideals and \(e_i \ge 1\).
\end{thm}

\begin{proof}
\textbf{Existence:} By Lemma \ref{product of non-zero prime ideals}, \(I\) contains a product of prime ideals. Choose such a product minimal with respect to inclusion. We prove equality. Suppose \(I\) strictly contains \(\mathfrak{p}_1 \cdots \mathfrak{p}_r\). Since \(\mathfrak{p}_1\) is invertible (Lemma \ref{lem:inv}), multiply both sides by \(\mathfrak{p}_1^{-1}\). Then
\[
\mathfrak{p}_2 \cdots \mathfrak{p}_r \subseteq I \mathfrak{p}_1^{-1} \subseteq \mathcal{O}_K.
\]
Since \(I \mathfrak{p}_1^{-1}\) is an ideal and the product was minimal, equality must hold. Repeating gives \(I = \mathfrak{p}_1 \cdots \mathfrak{p}_r\). To get exponents, repeat the argument for powers.

\textbf{Uniqueness:} Suppose
\[
\mathfrak{p}_1 \cdots \mathfrak{p}_r = \mathfrak{q}_1 \cdots \mathfrak{q}_s.
\]
Then \(\mathfrak{p}_1 \mid \mathfrak{q}_1 \cdots \mathfrak{q}_s\). Since \(\mathfrak{p}_1\) is prime, \(\mathfrak{p}_1 = \mathfrak{q}_j\) for some \(j\). Multiply both sides by \(\mathfrak{p}_1^{-1}\) and repeat. This proves uniqueness.
\end{proof}


\section{Fractional Ideals}

\begin{defn}
A \textbf{fractional ideal} of \(\mathcal{O}_K\) is a non-zero \(\mathcal{O}_K\)-submodule \(I \subseteq K\) such that there exists \(0 \neq d \in \mathcal{O}_K\) with \(d I \subseteq \mathcal{O}_K\).
\end{defn}

\begin{cor}\label{cor:fractional ideal invertibility}
Every fractional ideal \(I\) is invertible, i.e., there exists a fractional ideal \(I^{-1} \subseteq K\) such that \(I I^{-1} = \mathcal{O}_K\).
\end{cor}

\begin{proof}
Choose \(d \in \mathcal{O}_K\) with \(d I \subseteq \mathcal{O}_K\). Let \(J = d I\), an ideal of \(\mathcal{O}_K\). By Theorem \ref{thm:to prime ideals}, \(J = \prod_{i=1}^r \mathfrak{p}_i^{e_i}\). Define
\[
I^{-1} := d \prod \mathfrak{p}_i^{-e_i},
\]
where \(\mathfrak{p}^{-e} := (\mathfrak{p}^{-1})^e\). Then
\[
I I^{-1} = \left( \frac{1}{d} J \right) \left( d \prod \mathfrak{p}_i^{-e_i} \right)
= J \prod \mathfrak{p}_i^{-e_i}
= \mathcal{O}_K.
\]
Thus every fractional ideal is invertible.
\end{proof}

\section{Serre's Lemma }

\begin{lemma}[Serre]\label{lem:Serre}
Let \(\alpha_1, \dots, \alpha_m \in K\) not all zero. For any proper ideal \(I \subset \mathcal{O}_K\), there exists \(\beta \in K\) such that:
\begin{enumerate}
    \item \(\beta \alpha_i \in \mathcal{O}_K\) for all \(i\);
    \item at least one \(\beta \alpha_i \notin I\).
\end{enumerate}
\end{lemma}
\begin{proof}
Let \(A = \sum_{i=1}^m \mathcal{O}_K \alpha_i\). This is a fractional ideal. By Corollary \ref{cor:fractional ideal invertibility}, \(A\) has an inverse \(A^{-1}\), so \(A A^{-1} = \mathcal{O}_K\).

Suppose for contradiction that for every \(\gamma \in A^{-1}\), we have \(\gamma A \subseteq I\). Then \(A A^{-1} \subseteq I\), hence \(\mathcal{O}_K \subseteq I\). Since \(I\) is a proper ideal of \(\mathcal{O}_K\), this is impossible. Therefore there exists \(\gamma \in A^{-1}\) such that \(\gamma A \not\subseteq I\).

Set \(\beta = \gamma\). Then for each \(i\),
\[
\beta \alpha_i \in \beta A \subseteq A^{-1} A = \mathcal{O}_K.
\]
Moreover, since \(\beta A \not\subseteq I\), at least one \(\beta \alpha_i \notin I\). This proves the lemma.
\end{proof}

\begin{cor}
Let \(c_1, \dots, c_m \in \mathbb{C}\) be algebraic over \(\mathbb{Q}\), not all zero. Let \(\mathfrak{m} \subset \mathcal{R}\) be a maximal ideal. Then there exists a non-zero algebraic integer \(\lambda\) such that:
\begin{enumerate}
    \item \(\lambda c_i \in \mathcal{R}\) for all \(i\);
    \item at least one \(\lambda c_i \notin \mathfrak{m}\).
\end{enumerate}
\end{cor}

\begin{proof}
Let \(K = \mathbb{Q}(c_1, \dots, c_m)\), and let \(\mathcal{O}_K = K \cap \mathcal{R}\). Since \(c_i \in K\), we apply Lemma \ref{lem:Serre} with \(I = \mathfrak{m} \cap \mathcal{O}_K\). This gives \(\beta \in K\) such that \(\beta c_i \in \mathcal{O}_K\) for all \(i\), and at least one \(\beta c_i \notin I\). Then \(\lambda = \beta\) is an algebraic integer, and the conclusion follows.
\end{proof}


\section{Decomposition and Inertia Groups}

Let \(L/K\) be a finite Galois extension of number fields, and set \(G = \mathrm{Gal}(L/K)\). Let \(\mathfrak{p}\) be a nonzero prime ideal of \(\mathcal{O}_K\), and let \(\mathfrak{P}\) be a prime ideal of \(\mathcal{O}_L\) lying over \(\mathfrak{p}\), i.e. \(\mathfrak{P} \cap \mathcal{O}_K = \mathfrak{p}\).

\begin{defn}[Decomposition Group]
The \emph{decomposition group} of \(\mathfrak{P}\) is the subgroup of \(G\) consisting of all automorphisms that preserve the ideal \(\mathfrak{P}\):
\[
D(\mathfrak{P}) := \{ \sigma \in G : \sigma(\mathfrak{P}) = \mathfrak{P} \}.
\]
\end{defn}

\begin{defn}[Inertia Group]
The \emph{inertia group} of \(\mathfrak{P}\) is the kernel of the natural restriction map from \(D(\mathfrak{P})\) to the automorphism group of the residue field extension:
\[
I(\mathfrak{P}) := \{ \sigma \in D(\mathfrak{P}) : \sigma(x) \equiv x \pmod{\mathfrak{P}} \text{ for all } x \in \mathcal{O}_L \}.
\]
Equivalently, \(I(\mathfrak{P})\) consists of those automorphisms in \(D(\mathfrak{P})\) that act trivially on the residue field \(\mathcal{O}_L/\mathfrak{P}\).
\end{defn}

\begin{remark}
Both \(D(\mathfrak{P})\) and \(I(\mathfrak{P})\) are subgroups of \(G\). Moreover, \(I(\mathfrak{P})\) is a normal subgroup of \(D(\mathfrak{P})\), since it is the kernel of a group homomorphism.
\end{remark}

We now recall the standard invariants attached to the pair \(\mathfrak{P}/\mathfrak{p}\).

\begin{defn}[Ramification Index and Residue Degree]
Let \(\mathfrak{p}\mathcal{O}_L = \prod_{i=1}^{g} \mathfrak{P}_i^{e_i}\) be the factorization of \(\mathfrak{p}\) in \(\mathcal{O}_L\). The integer \(e_i\) is the \emph{ramification index} of \(\mathfrak{P}_i\) over \(\mathfrak{p}\). The residue field extension
\[
(\mathcal{O}_L/\mathfrak{P}_i) / (\mathcal{O}_K/\mathfrak{p})
\]
is a finite extension of finite fields; its degree is the \emph{residue degree} \(f_i := [\mathcal{O}_L/\mathfrak{P}_i : \mathcal{O}_K/\mathfrak{p}]\).
\end{defn}

\begin{thm}[Fundamental Identity]
With notation as above, for any \(\mathfrak{P} \mid \mathfrak{p}\),
\[
[L:K] = e(\mathfrak{P}/\mathfrak{p}) \, f(\mathfrak{P}/\mathfrak{p}) \, g,
\]
where \(g\) is the number of distinct prime ideals of \(\mathcal{O}_L\) lying above \(\mathfrak{p}\).
\end{thm}

\begin{proof}[Justification]
The Galois group \(G\) acts transitively on the set \(\{\mathfrak{P}_1, \dots, \mathfrak{P}_g\}\). By the orbit-stabilizer theorem,
\[
|G| = |D(\mathfrak{P})| \cdot g.
\]
A standard result in the theory of Dedekind domains (proved by comparing the \(\mathfrak{p}\)-adic valuations and the orders of the residue fields) gives
\[
|D(\mathfrak{P})| = e(\mathfrak{P}/\mathfrak{p}) \, f(\mathfrak{P}/\mathfrak{p}).
\]
Since \(|G| = [L:K]\), the identity follows.
\end{proof}

\begin{cor}
The restriction homomorphism
\[
\varphi: D(\mathfrak{P}) \longrightarrow \mathrm{Gal}((\mathcal{O}_L/\mathfrak{P}) / (\mathcal{O}_K/\mathfrak{p}))
\]
has kernel \(I(\mathfrak{P})\). In particular,
\[
|D(\mathfrak{P})| = |I(\mathfrak{P})| \cdot f(\mathfrak{P}/\mathfrak{p}).
\]
Since \(|D(\mathfrak{P})| = e f\), it follows that
\[
|I(\mathfrak{P})| = e(\mathfrak{P}/\mathfrak{p}).
\]
\end{cor}

\begin{proof}
The restriction map
\[
\varphi: D(\mathfrak{P}) \longrightarrow \mathrm{Gal}\big((\mathcal{O}_L/\mathfrak{P}) / (\mathcal{O}_K/\mathfrak{p})\big)
\]
is surjective. This is a standard fact in algebraic number theory (see, for example, Neukirch, \emph{Algebraic Number Theory}, Chapter I, Proposition 9.6). Briefly, one uses that the residue field extension is finite and separable, and that the Galois group acts transitively on the prime ideals above \(\mathfrak{p}\). The induced map on the decomposition group surjects onto the Galois group of the residue field extension.

By the first isomorphism theorem,
\[
D(\mathfrak{P})/I(\mathfrak{P}) \cong \mathrm{Gal}\big((\mathcal{O}_L/\mathfrak{P}) / (\mathcal{O}_K/\mathfrak{p})\big).
\]
The residue field extension has degree \(f(\mathfrak{P}/\mathfrak{p})\), so
\[
|D(\mathfrak{P})| = |I(\mathfrak{P})| \cdot f(\mathfrak{P}/\mathfrak{p}).
\]
On the other hand, the fundamental identity gives
\[
|D(\mathfrak{P})| = e(\mathfrak{P}/\mathfrak{p}) \, f(\mathfrak{P}/\mathfrak{p}).
\]
Comparing the two equations yields
\[
|I(\mathfrak{P})| = e(\mathfrak{P}/\mathfrak{p}).
\]
\end{proof}

\begin{defn}[Unramified Prime]
We say that \(\mathfrak{p}\) is \emph{unramified} in \(L/K\) if \(e(\mathfrak{P}/\mathfrak{p}) = 1\) for all (equivalently, for one) \(\mathfrak{P} \mid \mathfrak{p}\). By the corollary, this is equivalent to
\[
I(\mathfrak{P}) = \{1\}.
\]
\end{defn}

\section{The Frobenius Element}

Assume now that \(\mathfrak{p}\) is unramified in \(L/K\). Then \(I(\mathfrak{P}) = \{1\}\), so the restriction map
\[
\varphi: D(\mathfrak{P}) \longrightarrow \mathrm{Gal}((\mathcal{O}_L/\mathfrak{P}) / (\mathcal{O}_K/\mathfrak{p}))
\]
is an isomorphism. The codomain is the Galois group of a finite extension of finite fields.

Let \(q := N(\mathfrak{p}) = |\mathcal{O}_K/\mathfrak{p}|\). By the theory of finite fields, the Galois group \(\mathrm{Gal}(\mathbb{F}_{q^f}/\mathbb{F}_q)\) is cyclic of order \(f\), generated by the Frobenius automorphism
\[
\Phi: \mathbb{F}_{q^f} \to \mathbb{F}_{q^f}, \qquad \Phi(x) = x^{q}.
\]

\begin{defn}[Frobenius Element at \(\mathfrak{P}\)]
Let \(\sigma \in D(\mathfrak{P})\) be the unique element such that \(\varphi(\sigma) = \Phi\). That is, \(\sigma\) is characterized by the two properties:
\[
\sigma(\mathfrak{P}) = \mathfrak{P}, \qquad
\sigma(x) \equiv x^{N(\mathfrak{p})} \pmod{\mathfrak{P}} \quad \text{for all } x \in \mathcal{O}_L.
\]
We denote this element by
\[
\mathrm{Frob}_{\mathfrak{P}} := \sigma.
\]
\end{defn}

\begin{prop}[Independence of the Choice of \(\mathfrak{P}\)]
Let \(\mathfrak{P}'\) be another prime ideal of \(\mathcal{O}_L\) lying over \(\mathfrak{p}\). Then there exists \(g \in G\) such that \(\mathfrak{P}' = g(\mathfrak{P})\). Moreover,
\[
\mathrm{Frob}_{\mathfrak{P}'} = g \,\mathrm{Frob}_{\mathfrak{P}}\, g^{-1}.
\]
\end{prop}

\begin{proof}
Since \(G\) acts transitively on the primes above \(\mathfrak{p}\), choose \(g \in G\) with \(g(\mathfrak{P}) = \mathfrak{P}'\). Let \(\sigma = \mathrm{Frob}_{\mathfrak{P}}\). For any \(x \in \mathcal{O}_L\),
\[
(g\sigma g^{-1})(\mathfrak{P}') = g\sigma(\mathfrak{P}) = g(\mathfrak{P}) = \mathfrak{P}',
\]
and
\[
(g\sigma g^{-1})(x) \equiv x^{N(\mathfrak{p})} \pmod{\mathfrak{P}'}.
\]
Thus \(g\sigma g^{-1}\) satisfies the characterizing properties of \(\mathrm{Frob}_{\mathfrak{P}'}\). By uniqueness, the equality holds.
\end{proof}

\begin{defn}[Frobenius Conjugacy Class]
The conjugacy class of \(\mathrm{Frob}_{\mathfrak{P}}\) in \(G\) depends only on the underlying prime \(\mathfrak{p}\) of \(K\). We denote this conjugacy class by
\[
\mathrm{Frob}_{\mathfrak{p}} \subset G.
\]
It is called the \emph{Frobenius conjugacy class at \(\mathfrak{p}\)}.
\end{defn}

\section{Local Factors of the Artin L-Function}

Let \(\rho: G \to GL(V)\) be a finite-dimensional complex representation of \(G = \mathrm{Gal}(L/K)\). We now define the local factor of the Artin L-function attached to \(\rho\) at each prime \(\mathfrak{p}\) of \(K\).

\subsection{The Unramified Case}

If \(\mathfrak{p}\) is unramified in \(L/K\), the Frobenius conjugacy class \(\mathrm{Frob}_{\mathfrak{p}}\) is well-defined. The local factor is given by
\[
L_{\mathfrak{p}}(s, \rho) := \det\!\left( I - \rho(\sigma) N(\mathfrak{p})^{-s} \right)^{-1},
\]
where \(\sigma \in \mathrm{Frob}_{\mathfrak{p}}\) is any representative.

\begin{remark}
This definition is independent of the choice of \(\sigma\), because if \(\sigma' = g\sigma g^{-1}\), then
\[
\det(I - \rho(\sigma') N^{-s}) = \det(I - \rho(g)\rho(\sigma)\rho(g)^{-1} N^{-s}) = \det(I - \rho(\sigma) N^{-s}),
\]
since determinant is invariant under conjugation.
\end{remark}

\subsection{The Ramified Case}

If \(\mathfrak{p}\) is ramified (\(e(\mathfrak{P}/\mathfrak{p}) > 1\)), the Frobenius element is not well-defined as a conjugacy class. In this case, we fix a prime \(\mathfrak{P} \mid \mathfrak{p}\). The quotient group
\[
D(\mathfrak{P}) / I(\mathfrak{P})
\]
is isomorphic to the Galois group of the residue field extension, which is cyclic and generated by the Frobenius automorphism. Choose any lift \(\sigma \in D(\mathfrak{P})\) whose image in \(D(\mathfrak{P})/I(\mathfrak{P})\) corresponds to the Frobenius automorphism of the residue fields.

Consider the subspace of \(V\) fixed by the inertia group:
\[
V^{I(\mathfrak{P})} := \{ v \in V : \rho(\tau)v = v \text{ for all } \tau \in I(\mathfrak{P}) \}.
\]
This subspace is preserved by \(\rho(\sigma)\), because \(I(\mathfrak{P})\) is normal in \(D(\mathfrak{P})\). The local factor at the ramified prime \(\mathfrak{p}\) is defined as
\[
L_{\mathfrak{p}}(s, \rho) := \det\!\left( I - \rho(\sigma) N(\mathfrak{p})^{-s} \mid V^{I(\mathfrak{P})} \right)^{-1},
\]
where the determinant is taken on the invariant subspace \(V^{I(\mathfrak{P})}\).

\begin{prop}[Well-Definedness of the Ramified Local Factor]
The ramified local factor \(L_{\mathfrak{p}}(s, \rho)\) does not depend on the choices of:
\begin{enumerate}
    \item the prime \(\mathfrak{P} \mid \mathfrak{p}\),
    \item the lift \(\sigma \in D(\mathfrak{P})\) of the Frobenius element.
\end{enumerate}
\end{prop}

\begin{proof}
(1) If \(\mathfrak{P}' = g(\mathfrak{P})\), then \(I(\mathfrak{P}') = g I(\mathfrak{P}) g^{-1}\), and \(V^{I(\mathfrak{P}')} = \rho(g) V^{I(\mathfrak{P})}\). The determinant is invariant under this conjugation.

(2) If \(\sigma_1, \sigma_2\) are two lifts of the same Frobenius element in \(D(\mathfrak{P})/I(\mathfrak{P})\), then \(\sigma_1 \sigma_2^{-1} \in I(\mathfrak{P})\). Hence \(\rho(\sigma_1)\) and \(\rho(\sigma_2)\) differ by an element of \(\rho(I(\mathfrak{P}))\), which acts trivially on \(V^{I(\mathfrak{P})}\). Thus their restrictions to \(V^{I(\mathfrak{P})}\) are identical, and so are the local factors.
\end{proof}

\subsection{The Global Artin L-Function}

The global Artin L-function attached to \(\rho\) is defined as the Euler product over all prime ideals \(\mathfrak{p}\) of \(\mathcal{O}_K\):
\[
L(s, \rho, L/K) := \prod_{\mathfrak{p}} L_{\mathfrak{p}}(s, \rho),
\]
where the product converges absolutely for \(\Re(s) > 1\). Using the local factors defined above, this can be written explicitly as:
\[
L(s, \rho, L/K) =
\prod_{\mathfrak{p} \text{ unramified}}
\det\!\left( I - \rho(\mathrm{Frob}_{\mathfrak{p}}) N(\mathfrak{p})^{-s} \right)^{-1}
\cdot
\prod_{\mathfrak{p} \text{ ramified}}
\det\!\left( I - \rho(\sigma_{\mathfrak{p}}) N(\mathfrak{p})^{-s} \mid V^{I(\mathfrak{P}_{\mathfrak{p}})} \right)^{-1}.
\]

\begin{remark}
For a one-dimensional representation \(\rho = \chi\), the determinant is a scalar. In the unramified case, this simplifies to
\[
L_{\mathfrak{p}}(s, \chi) = \left(1 - \chi(\mathrm{Frob}_{\mathfrak{p}}) N(\mathfrak{p})^{-s}\right)^{-1},
\]
which is the standard Euler factor for Hecke L-functions.
\end{remark}